% ============================================================
% Section 2.2 & 2.4: New Related Work sections for WISE
% Quality Assurance and Traceability in RAG systems
% ============================================================

\subsection{Quality Assurance and Verification in RAG}

Retrieval-Augmented Generation mitigates LLM hallucination by grounding outputs
in retrieved evidence, but early RAG systems provided no mechanism to verify
that the generated answer is actually supported by the retrieved passages. This
verification gap has motivated recent work on answer quality assessment and
factual consistency checking.

Self-RAG~\cite{asai2023selfrag} introduces reflection tokens that critique
retrieval quality and answer faithfulness at the token level during generation,
enabling the model to decide when to retrieve and when to self-correct. This
approach achieves strong gains on knowledge-intensive tasks but requires
parameter training via reinforcement learning, limiting its applicability to
frontier LLMs accessible only via API. Moreover, the critique operates at the
token level rather than the retrieval-component level, making it difficult to
isolate \emph{which} retrieval step failed.

FAIR-RAG~\cite{fairrag2025} addresses evidence gaps by prompting the LLM to
assess whether retrieved passages contain sufficient information to answer the
question, and triggers re-retrieval when gaps are detected. This provides a
training-free quality check but operates at the evidence-set level rather than
the sub-query level, so when re-retrieval is triggered, the system must
re-retrieve the full evidence set rather than targeting the specific failing
component. FAIR-RAG also does not produce a structured diagnostic record
identifying \emph{which} fact is missing or \emph{why} retrieval failed, limiting
its auditability for Web systems requiring traceable reasoning.

CoRAG~\cite{lyu2025corag} iteratively refines retrieval through chain-of-thought
planning: the LLM decides whether to retrieve again based on answer confidence,
and restarts the full pipeline when refinement is needed. This provides
adaptivity but restarts from scratch rather than surgically repairing the failing
component, wasting computation on re-retrieving evidence that was already correct.

VERIFY~\cite{verify2024} and FactScore~\cite{factscore2023} decompose generated
answers into atomic claims and verify each claim against retrieved passages or
a knowledge source, enabling fine-grained factual consistency checking. These
systems excel at post-hoc verification but do not provide a retrieval-repair
mechanism: they can detect that an answer is wrong but cannot fix the underlying
retrieval failure that caused it.

The common thread across these systems is that none provides \emph{sub-query-level}
quality diagnostics combined with a surgical repair mechanism. Self-RAG and
FAIR-RAG detect quality issues but cannot pinpoint which retrieval component
failed; CoRAG restarts globally rather than repairing targeted components;
VERIFY and FactScore verify but do not repair. \sys fills this gap by verifying
factual consistency at the sub-query level, identifying the minimum failing unit,
and enabling surgical repair that preserves successful retrieval work while
revising only what failed.

\subsection{Traceability and Explainability in Web QA Systems}

Trustworthy Web information systems require not only correct answers but
\emph{traceable reasoning}---a structured record of how the system arrived at
its answer, which retrieval steps succeeded, and where failures occurred. This
traceability is essential for regulated domains (healthcare, finance, legal QA)
where answers must be auditable, and for production deployment where system
maintainers need diagnostic signals to debug failures.

Explainable QA has a rich history. Early neural QA systems~\cite{rajpurkar2016squad}
provided answer span predictions with attention weights, but attention-based
explanations are notoriously unreliable~\cite{jain2019attention}. Chain-of-thought
prompting~\cite{wei2022cot} improved interpretability by eliciting intermediate
reasoning steps, but the generated reasoning is post-hoc rationalization rather
than a faithful trace of the retrieval process, and recent work has shown that
CoT explanations often contradict the model's actual decision
process~\cite{turpin2023cot}.

For RAG systems, several approaches have emerged for tracing retrieval provenance.
Retrieve-and-Read systems~\cite{karpukhin2020dpr} return the top-$k$ retrieved
passages alongside the answer, providing evidence provenance but no record of
\emph{why} those passages were selected or which parts were used. IRCOT and
HopRAG interleave retrieval and reasoning, logging each retrieval step, but the
logs are unstructured and provide no diagnostic signal indicating which step
failed or why.

GraphRAG systems~\cite{edge2024graphrag,guo2024lightrag} anchor answers to
specific KG triples, providing structured evidence provenance superior to
passage-level attribution. However, none of the existing GraphRAG systems
produces a diagnostic record identifying which entity-linking step failed, which
bridge entity was missed, or why a sub-query could not be resolved. When a
GraphRAG system returns a wrong answer, the user sees the incorrect output and
the retrieved triples but has no structured explanation of where the reasoning
broke down.

A-RAG~\cite{arag2026} logs the full trajectory of tool calls (keyword search,
semantic search, chunk read), providing detailed provenance of the retrieval
process. This improves transparency but remains action-trace logging rather than
diagnostic output: the log records \emph{what} the system did but not \emph{why}
a retrieval step failed or \emph{what} should be revised to fix it.

\sys addresses this gap by producing \emph{structured diagnostic records} at the
sub-query level. When the VA detects a quality failure, it emits a natural-language
diagnostic string identifying (1)~which sub-query failed, (2)~the diagnosed
failure type (entity not found, incomplete evidence, contradictory triples), and
(3)~a concrete revision suggestion. This structured output serves dual purposes:
it enables surgical repair (Section~\ref{sec:quality_analysis}) and it provides
traceable reasoning for auditability. Every answer comes with a record of which
sub-queries were resolved successfully, which required revision, and why---making
\sys the first GraphRAG system whose reasoning process is fully auditable.

The distinction is architectural: prior systems treat traceability as logging
(record what happened), while \sys treats traceability as \emph{diagnostic output}
(identify what failed and why). This shift from passive logging to active
diagnostics is what enables both surgical repair and trustworthy deployment in
Web systems serving critical applications.

